\begin{table}[htbp]
\centering
\caption{Price clustering by market-capitalisation quintile}
\label{tab:size}
\small
\begin{threeparttable}
\begin{tabular}{@{}lrrrrrr@{}}
\toprule
Quintile & $M^{0}$ (\%) & $M^{BLarge0}$ (\%) & SE & $M^{BLarge0}$, \textyen 1 grid (\%) & 0.1-tick share (\%) & Stock-days \\
\midrule
Q1 & 14.45 & 14.83 & (0.304) & 14.83 & 0.0 & 15,743 \\
Q2 & 14.66 & 15.55 & (0.291) & 15.55 & 0.0 & 15,689 \\
Q3 & 14.20 & 14.77 & (0.257) & 14.77 & 0.0 & 15,683 \\
Q4 & 13.12 & 13.81 & (0.220) & 13.70 & 1.9 & 15,450 \\
Q5 & 14.70 & 15.82 & (0.250) & 15.51 & 12.2 & 14,764 \\
\bottomrule
\end{tabular}
\begin{tablenotes}[flushleft]\footnotesize
\item Quintiles are formed once per stock on its median market capitalisation over the sample, so a stock does not drift between quintiles as its price moves. Q1 is the smallest. The pooled columns mix tick regimes; the \textyen 1-grid column removes the mechanical elevation that 0.1-yen days carry, and the 0.1-tick share column shows how unevenly those days fall across quintiles.
\end{tablenotes}
\end{threeparttable}
\end{table}
